\subsection{媒介変数表示}

円や,縦向きの接線を持つ曲線は,1つの関数 $y=f(x)$ のグラフとして扱いにくい.そこで
\[
x=x(t),\qquad y=y(t),\qquad \alpha\leq t\leq\beta
\]
と表された曲線の長さを考えよう.\,$x(t),y(t)$ は連続微分可能であるとする.

$t$ を $\Delta t$ だけ増やしたとき,横の変化はほぼ $x'(t)\Delta t$,縦の変化はほぼ $y'(t)\Delta t$ である.よって,対応する線分の長さはほぼ
\[
\sqrt{\{x'(t)\Delta t\}^2+\{y'(t)\Delta t\}^2}
 =\sqrt{x'(t)^2+y'(t)^2}\,\Delta t
\]
となる.これを足し合わせた折れ線の極限として
\begin{equation}\label{式:媒介曲線の弧長}
L=\int_\alpha^\beta\sqrt{x'(t)^2+y'(t)^2}\dd t
\end{equation}
を得る.導関数が連続であることから近似の誤差は区間全体で小さくできるが,ここではその評価の詳細は省き,この式を用いる.途中で $\sqrt{x'(t)^2+y'(t)^2}$ が $0$ になる点があっても,この式はそのまま使える.

ここで1つ断りをいれておこう.この積分は,\,$t$ に沿って点が進んだ長さを数える.同じ部分を2度たどれば2度数える.たとえば $x=\cos t$,\,$y=\sin t$ を $0\leq t\leq4\pi$ で積分すると $4\pi$ になるが,図形としての単位円の周の長さは $2\pi$ である.図形そのものの長さを求めるなら,1度だけたどる区間を選ぶ.

媒介変数表示は面積にも使える.\,$x(t)$ が単調増加し,\,$y(t)\geq0$ であるとする.このとき曲線は $x$ の関数のグラフであり,曲線と $x$ 軸の間の面積は $\int_{x(\alpha)}^{x(\beta)}y\dd x$ である.\,$x=x(t)$ と置換すると $\dd x=x'(t)\dd t$ なので,
\begin{equation}\label{式:媒介曲線の面積}
S=\int_\alpha^\beta y(t)x'(t)\dd t
\end{equation}
となる.\,$x(t)$ が途中で折り返すなら,区間を分けて,曲線が囲む部分を確認する.この式に絶対値を付けるだけで,いつでも図形の面積になるわけではない.縦の帯が重なれば,同じ面積を2度数えるからである.

\begin{bookpoint}
\textbf{POINT}\par
長さは $\sqrt{x'(t)^2+y'(t)^2}$ を,面積は $y(t)x'(t)$ を $t$ で積分する.根号を外すときは $t$ の範囲で符号を確かめる.
\end{bookpoint}

\begin{example}{媒介曲線の面積と長さ}
$a>0$ とする.曲線
\[
x=a(t-\sin t),\qquad y=a(1-\cos t),\qquad 0\leq t\leq2\pi
\]
の長さと,この曲線と $x$ 軸に囲まれた部分の面積を求めよ.
\end{example}

\paragraphbf{方針}
長さは式\eqref{式:媒介曲線の弧長}で求め,根号の中を半角の公式で2乗の形にする.面積は $x(t)$ が単調増加であることを確かめてから式\eqref{式:媒介曲線の面積}を使う.

\paragraphbf{解答}
長さ $8a$,面積 $3\pi a^2$

\paragraphbf{解法}
$x'(t)=a(1-\cos t)$,\,$y'(t)=a\sin t$ である.

\textbf{長さ}\quad 根号の中は
\begin{align*}
x'(t)^2+y'(t)^2
 &=a^2\{(1-\cos t)^2+\sin^2t\}\\
 &=a^2(2-2\cos t)\\
 &=4a^2\sin^2\frac t2
\end{align*}
となる.平方根は $2a\left|\sin\frac t2\right|$ である.\,$0\leq t\leq2\pi$ では $0\leq\frac t2\leq\pi$ なので $\sin\frac t2\geq0$ であり,
\begin{align*}
L&=2a\int_0^{2\pi}\sin\frac t2\dd t\\
 &=\left[-4a\cos\frac t2\right]_0^{2\pi}\\
 &=4a-(-4a)\\
 &=8a.
\end{align*}

\textbf{面積}\quad $x'(t)=a(1-\cos t)$ は $0<t<2\pi$ で正なので,\,$x$ は区間全体で単調増加する.また $y\geq0$ で,両端では $y=0$ である.よって
\begin{align*}
S&=\int_0^{2\pi}a(1-\cos t)\cdot a(1-\cos t)\dd t\\
 &=a^2\int_0^{2\pi}(1-2\cos t+\cos^2t)\dd t\\
 &=a^2\int_0^{2\pi}\left(\frac32-2\cos t+\frac12\cos2t\right)\dd t\\
 &=a^2\cdot\frac32\cdot2\pi\\
 &=3\pi a^2.
\end{align*}
$\cos t$ と $\cos2t$ の項は,\,$0$ から $2\pi$ まで積分すると $0$ になる.
